\documentclass[tikz,border=5pt]{standalone}
\usepackage[UTF8]{ctex}
\usepackage{amsmath,amssymb}
\usetikzlibrary{arrows.meta,calc,angles,quotes,decorations.pathmorphing,patterns}
\definecolor{ink}{HTML}{243447}\definecolor{blue}{HTML}{2765A2}
\definecolor{green}{HTML}{3F7D5B}\definecolor{red}{HTML}{B74739}
\tikzset{body/.style={draw=ink,thick},vec/.style={draw=blue,thick,-{Latex}},guide/.style={gray!60,dashed},dim/.style={ink,{Latex}-{Latex}},line cap=round,line join=round}
\begin{document}

\begin{tikzpicture}
\node[draw=ink,rounded corners,minimum width=2.1cm,minimum height=.9cm,align=center](A)at(0,0){液态水\\雾化功率 $\eta P$};
\node[draw=ink,rounded corners,minimum width=2cm,minimum height=.9cm,align=center](B)at(3.3,0){雾滴\\半径 $r$};
\node[draw=ink,rounded corners,minimum width=2.5cm,minimum height=.9cm,align=center](C)at(6.9,0){水蒸气\\分子源 $\Gamma$};
\draw[vec](A)--node[above]{$\dot V_{\mathrm{atom}}$}(B);
\draw[vec](B)--node[above]{$\eta_2$}(C);
\draw[vec](C.east)--++(1.1,0)node[right]{孔隙逸出};
\node[below]at(4,-.85){雾滴与自由水分子分别统计};
\end{tikzpicture}
\begin{tikzpicture}
\fill[blue!6](-1.3,0)arc(180:0:1.3)--cycle;
\draw[body](-1.3,0)arc(180:0:1.3);
\draw[guide](-1.3,0)arc(180:0:1.3 and .35);
\draw[body](-1.3,0)arc(180:360:1.3 and .35);
\draw[dim](0,0)--node[below]{$r$}(1.3,0);
\draw[vec,red](-1.3,0)--(-1.3,-.7);
\draw[vec,red](1.3,0)--(1.3,-.7)node[right]{$\sigma\,d\ell$};
\draw[vec,red](0,-.35)--(0,-1.05);
\node[red,below]at(0,-1.1){周边合力 $2\pi r\sigma$};
\node[blue]at(0,.75){$P_1$};
\node at(1.9,1.1){外界 $P_0$};
\draw[vec](0,1.25)--(0,2.2)node[right]{$(P_1-P_0)\pi r^2$};
\node[below]at(0,-1.7){上半滴：一个界面，截口周长 $2\pi r$};
\end{tikzpicture}
\begin{tikzpicture}
\draw[body](0,0)--(0,1.4)--(1.4,1.4)--(1.4,0)--cycle;
\draw[body](0,1.4)--(.45,1.75)--(1.85,1.75)--(1.4,1.4);
\draw[body](1.4,0)--(1.85,.35)--(1.85,1.75);
\fill[blue!8](.02,.02)rectangle(1.38,1.1);
\node at(.7,.5){$V_0$};
\draw[vec](2.25,.75)--(3.45,.75);
\foreach\x in{4,4.6,5.2}\foreach\y in{.15,.75,1.35}{
 \fill[blue!8](\x,\y)circle(.21);
 \draw[body](\x,\y)circle(.21);
 \draw[guide](\x-.21,\y)arc(180:0:.21 and .065);
 \draw[gray!60](\x-.21,\y)arc(180:360:.21 and .065);
}
\node[above]at(4.6,1.9){$N_d$ 个半径 $r$ 的小球};
\node[below,align=center]at(2.8,-.6){$N_d(4\pi r^3/3)=V_0$\\$A_d=N_d4\pi r^2=3V_0/r$};
\end{tikzpicture}
\begin{tikzpicture}
\draw[body](0,0)rectangle(4.2,2.5);
\node[align=center]at(1.8,1.5){室内 $n(t)$\\体积 $V_R$};
\draw[vec]( -.9,.9)--(1,.9)node[above left]at(-.4,.9){$\Gamma$};
\draw[white,line width=4pt](4.2,.8)--(4.2,1.6);
\draw[body](4.2,.8)--(4.7,.8)--(4.7,1.6)--(4.2,1.6);
\node[above]at(4.45,1.7){$S$};
\draw[vec](3.9,1.2)--(5.9,1.2)node[right]{$nS\bar v/6$};
\node[anchor=west]at(5.6,2.3){室外近似干燥};
\begin{scope}[shift={(1.8,-1.4)},scale=.75]
\foreach\p in{(1.1,0),(-1.1,0),(0,1),(0,-1),(.65,.5),(-.65,-.5)}
\draw[vec,green](0,0)--\p;
\node[green,right]at(1.2,0){六方向等概率};
\end{scope}
\node[below]at(2.9,-2.5){稳态：$\Gamma=n_{\mathrm{ss}}S\bar v/6$};
\end{tikzpicture}
\end{document}
